% Responsable(s): por asignar
% Objetivo: comparar cuantitativamente latencia, ancho de banda, tiempo de
% acceso y demás métricas reportadas por los artículos seleccionados.

% Redacción pendiente.
\subsubsection{HBM frente a DDR en la supercomputadora Aurora}

Intel y el Laboratorio Nacional Argonne pusieron a prueba, dentro de la supercomputadora Aurora,
los dos tipos de memoria que conviven en sus nodos de cómputo: la memoria de alto ancho de
banda (HBM) integrada en el procesador Xeon Max, y la DDR tradicional del mismo sistema. El
benchmark que usaron, STREAM Triad, midió el ancho de banda sostenido bajo distintas
combinaciones de clustering (Quad y SNC4) y modos de memoria (Flat y Cache). Con socket
completo en modo SNC4, HBM llegó a 840 GB/s, más del triple de los 245 GB/s que alcanzó
DDR \cite{ibeid2025aurora}. Pero la historia cambia con la latencia: al medirla sin carga usando Intel
Memory Latency Checker, DDR resultó más rápida en todas las configuraciones probadas, con
una ventaja de más del 10\% en modo SNC4. La conclusión de los autores es simple pero útil: ni
HBM ni DDR ganan en todo, así que la mejor opción depende de si la aplicación necesita más
ancho de banda o menos latencia.

\subsubsection{El marco Mess: una vista unificada del sistema de memoria}

Esmaili-Dokht y su equipo ---trabajo que quedó como finalista al mejor artículo en MICRO 2024
--- construyeron algo distinto a lo habitual: en vez de reportar un solo número de latencia o ancho
de banda, propusieron caracterizar la memoria con curvas completas de ancho de banda contra
latencia, a las que llamaron framework Mess. Corrieron su benchmark sobre servidores reales de
Intel, AMD, IBM, Fujitsu y Amazon, además de GPUs NVIDIA, cubriendo DDR4, DDR5, HBM2
y HBM2E. Un dato llamativo: dos servidores con la misma memoria DDR4 pueden tener latencias
muy distintas. Uno con arquitectura Cascade Lake midió 85 ns sin carga, mientras que otro con
Zen2 ---también DDR4, con temporizaciones casi idénticas--- llegó a casi 30 ns más
\cite{esmailidokht2024mess}. Esto deja claro que la latencia no depende solo del chip de memoria, sino de
toda la microarquitectura que lo rodea: caché, interconexión, número de núcleos. Además, algo
útil para el resto del grupo: los autores probaron que simuladores tan usados como gem5 no logran
replicar bien estos comportamientos reales, así que conviene tener cuidado si en algún momento
se citan resultados de simulación en vez de mediciones directas.

\subsubsection{Memoria heterogénea con enlaces de coherencia de caché: CXL, NVLink-C2C e Infinity Fabric}

Con la llegada de CXL, NVLink-C2C e Infinity Fabric, la memoria dejó de limitarse a los canales
DIMM de siempre, y eso es justo lo que investigó este grupo de UC San Diego con su banco de
pruebas Heimdall. Probaron diez tipos de servidores distintos combinando CPUs de varios
fabricantes con DDR5 DIMM, LPDDR5x, HBM3, HBM3e y memoria expandida vía CXL. Los
sistemas DDR5 de cuatro canales saturan su ancho de banda cerca de 107 GiB/s; con ocho canales,
la cifra sube considerablemente antes de que la latencia empiece a dispararse
\cite{wang2024hitchhiker}. Lo más interesante para esta comparación es el costo de usar CXL: un acceso remoto vía este
protocolo puede ser entre 50\% y 300\% más lento que un acceso local a DIMM. En números
concretos, mientras la DRAM local ronda los 100 ns, un dispositivo CXL basado en ASIC se
mueve entre 200 y 300 ns, y los prototipos FPGA llegan hasta 400 ns. Este artículo aporta justo lo
que le faltaba a los otros dos: tecnologías de memoria de última generación ---HBM3, HBM3e,
CXL--- con números concretos de cuánto cuesta, en latencia, expandir la memoria más allá de lo
tradicional.

\subsubsection{Controlador de memoria híbrido con eDRAM y MRAM para Big Data}

El cuarto artículo cambia de enfoque: en vez de medir memoria principal a nivel de sistema
completo, se mete en el diseño del controlador de memoria mismo. Fue presentado en SARTECO
2023, el congreso anual de la Sociedad Española de Arquitectura y Tecnología de Computadores,
por un equipo de la Universitat Politècnica de València. El problema que atacan es conocido: la
memoria no volátil (NVRAM) ofrece mucha más capacidad que la DRAM, pero paga esa
capacidad con latencias de acceso mayores. La solución que proponen es añadir una segunda
caché, más densa que la SRAM habitual, construida con eDRAM o MRAM, que absorbe los fallos
antes de que lleguen a NVRAM. Simulando cargas de trabajo típicas de Big Data, el esquema con
eDRAM bajó la penalización por fallo hasta un 50\%, y el de MRAM hasta un 80\%; a nivel de
sistema completo eso se tradujo en mejoras de rendimiento de 15\% y 23\% respectivamente
\cite{lurbe2023controlador}. Sirve porque es el único de los cuatro que no habla de una tecnología de memoria ya
establecida en el mercado, sino de cómo diseñar la jerarquía de caché para tecnologías que todavía
están emergiendo.

\subsubsection{Comparación cuantitativa}

La siguiente tabla resume los valores numéricos más relevantes reportados en los cuatro artículos,
organizados por tecnología y tipo de medición.

\begin{table}[ht]
\centering
\small
\begin{tabular}{p{4.2cm} p{3cm} p{3cm} p{3.2cm}}
\toprule
\textbf{Tecnología / configuración} &
\textbf{Ancho de banda medido} &
\textbf{Latencia medida} &
\textbf{Fuente} \\
\midrule

HBM (Xeon Max, modo SNC4) &
840 GB/s &
$\sim$121 ns (idle) &
Ibeid et al. (2025) \\

DDR (Xeon Max, modo SNC4) &
245 GB/s &
$\sim$96 ns (idle) &
Ibeid et al. (2025) \\

DDR4 (servidor Cascade Lake) &
--- &
85 ns (sin carga) &
Esmaili-Dokht et al. (2024) \\

DDR5 (Graviton 3) &
--- &
129 ns (sin carga) &
Esmaili-Dokht et al. (2024) \\

HBM2 (A64FX) &
--- &
122 ns (sin carga) &
Esmaili-Dokht et al. (2024) \\

DDR5 DIMM (4 canales, punto de saturación) &
107 GiB/s &
--- &
Wang et al. (2024) \\

Memoria CXL remota (ASIC) vs. DIMM local &
--- &
200--300 ns vs. $\sim$100 ns &
Wang et al. (2024) \\

Caché híbrida eDRAM/MRAM vs. SRAM sola &
--- &
$-50\%$ a $-80\%$ en penalización por fallo &
Lurbe et al. (2023) \\

\bottomrule
\end{tabular}
\caption{Comparación cuantitativa de tecnologías y configuraciones de memoria.}
\label{tab:comparacion-memoria}
\end{table}

Viendo estos números juntos, se repite una misma idea en los cuatro artículos: no existe una
memoria que gane en ancho de banda y en latencia al mismo tiempo. HBM y sus variantes se
llevan la ventaja en ancho de banda ---algo clave para cargas de trabajo muy paralelas, como HPC
o entrenamiento de modelos de IA--- pero pierden entre 10\% y 20\% de latencia frente a DDR
\cite{ibeid2025aurora,esmailidokht2024mess}. Salir de los canales DIMM tradicionales también
tiene un precio, sea por CXL o por una jerarquía de caché adicional: CXL suma entre 100 ns y 300
ns sobre un acceso local \cite{wang2024hitchhiker}, mientras que las cachés híbridas con eDRAM o
MRAM hacen justo lo contrario, intentan esconder la latencia de la NVRAM a cambio de un
controlador de memoria más complejo \cite{lurbe2023controlador}. Y algo que vale la pena recordar: la
latencia real que se mide en un sistema depende tanto de la memoria como de todo lo que la rodea
---caché, interconexión, cantidad de núcleos---, razón por la cual dos servidores con la misma
DDR4 pueden diferir en casi 30 ns \cite{esmailidokht2024mess}.